\section{Conclusion}
\label{sec:conclusion}
This paper revisits \cite{justinianoInvestmentShocksBusiness2010}, who claims that the investment specific technology is the main driving force of the business cycle, by testing an alternative hypothesis---the investment demand shock proposed by \cite{nireiOriginsPropagationAnimal}. Building a DSGE model with both the IST shock and the investment demand shock, the estimation results show that the investment demand shock can be another competing shock driving the business cycle. In the baseline result, the investment demand shock alone contributes about 40\% of the forecast error of output, and about 50\% of investment.

However, there are a few limitations in this paper. First, there is a lack of data to directly identify the investment demand shock, and thus the identification of the investment demand shock and the IST shock can only be based on their different effects on other state variables. Moreover, it is hard to test the argument across all combinations of parameter values, as the parameter dimension is high. Section~\ref{sec:results:inf} and \ref{sec:results:con} can only provide partial evidence on this argument.
Second, the baseline model cannot fit the data very well, and the posteriors of some parameters take extreme values, which is discussed in Section \ref{sec:estimation:dist} and \ref{sec:results:moments}. Third, the estimation results depend on the model specification. Expanding the model (Section~\ref{sec:robustness:sw}) or using an alternative stochastic process (Section~\ref{sec:robustness:iid}) can alter the result quite a lot.

Nevertheless, this paper highlights the insights provided by structural estimation, but also the fragility of such an approach. More macroeconomic effects of investment-side shocks can be discussed and compared, but the model specification should be selected carefully.